\section{Sarsa}
    Now, we introduce the most basic TD-algorithm, the \textbf{Sarsa}. Its name originates from the
    tuple $(s,a,r,s',a')$ it uses.
    
    In the last section, the TD algorithm of learning state value was mainly discussed. Here, we directly learn the
    action values instead of state values, since it can be combined with a policy improvement step to learn optimal policies.
    \subsection{Main Idea}
        Suppose we have trajectory generated from policy $\pi$: $(s_0,a_0,r_1,s_1,a_1,\dots,s_t,a_t,r_{t+1},s+{t+1},a_{t+1},\dots).$
        Note that the trajectory does not have to end. We use following TD-update every step for estimating the action values:
        \begin{align}
            q_{t+1}(s_t,a_t) &= q_t(s_t,a_t) + \alpha_t(s_t,a_t)\Big[ r_{t+1} + \gamma q_t(s_{t+1},a_{t+1}) - q_t(s_t,a_t) \Big] \label{eq:SarsaUpdate}\\
            q_{t+1}(s,a) &= q_t(s,a),\text{ for all }(s,a) \neq (s_t,a_t)
        \end{align}    

        In fact, the Sarsa solves following Bellman equation in a stochastic approximation manner.
        \begin{equation}
            q_\pi(s,a) = \E[R + \gamma q_\pi(S',A')|s,a] \text{ for all }(s, a) \label{eq:sarsaBellman}
        \end{equation}
        Equation \ref{eq:sarsaBellman} is the Bellman equation about action values. The proof is given below.
        \begin{proof}
            Note that
            \begin{align}
                p(s',a'|s,a) &= p(s'|s,a)p(a'|s',s,a) \\
                             &= p(s'|s,a)p(a'|s') \text{ (due to the Markov property)} \\
                             &= p(s'|s,a)\pi(a'|s').
            \end{align}
            Using this, the equation of Theorem \ref{thm:BellmanEq-AV} can be written as
            \begin{equation}
                q_\pi(s,a) = \sum_{r}rp(r|s,a) + \gamma \sum_{s',a'}q_\pi(s',a')p(s',a'|s,a)
            \end{equation}
        \end{proof}
    \subsection{Pseudocode}
        \Algorithm{sarsa}
        {Sarsa}
        {\inputitem{$\alpha_t(s,a) \in (0,1]$}{step size parameter satisfying RM condition}}
        {\inputitem{$q_*$}{the optimal policy}}
        {
            \tcp{initialize}
            \ForEach{$(s,a) \in \mathcal{S} \times \mathcal{A}$}
            {$q(s,a) \leftarrow
            \begin{cases}
                \text{arbitrary,} &\text{ if $s \in \mathcal{S}^+$ and $a \in \mathcal{A}(s)$} \\
                0 &\text{ if $s \in \mathcal{S} \setminus \mathcal{S}^+$}
            \end{cases}$\;}
            $\pi \leftarrow$ policy derived from $q$ (e.g. $\varepsilon$-greedy, softmax, GLIE)\;

            \ForEach{episodes}
            {
                initialize $s$\;
                choose action $a$ from $s$ using policy $\pi$\;
                \Repeat{s is in terminal state}
                {
                    take action $a$, observe $r$, $s'$\;
                    choose action $a'$ from $s'$ using policy derived from $\pi$\;
                    \tcp{update q value}
                    $q(s,a) \leftarrow q(s,a) + \alpha_t(s,a) \big[ r + \gamma q(s',a') - q(s,a) \big]$\;
                    \tcp{update policy}
                    update $\pi$ according to $q(s,a)$
                    $s \leftarrow s'$\;
                    $a \leftarrow a'$\;
                }
            }
            \Return{$q \simeq q_*$}
        }
    \subsection{Implementation Detail}
        I usually intialize q value without minding if the state is terminal or not. While traning the agent, when the agent meets terminal state,
        I just check it in place and set the $q_t(s_{t+1},a_{t+1}) = 0$. Actually I don't even change the q value of terminal state to 0. I just leave it untouched.
    \subsection{Convergence}
        \begin{thm}{Convergence of Sarsa}{convSarsa}
            Given a policy $\pi$, by the Sarsa algorithm \ref{eq:SarsaUpdate}, $q_t(s,a)$ converges almost surely to the action value $q_\pi(s,a)$ as $t \to \infty$ for all
            $(s,a)$ if $\sum_t \alpha_t(s,a) = \infty$ and $\sum_t \alpha_t^2(s,a) < \infty$ for all $(s,a)$.
        \end{thm}
        \begin{proof}
            TODO: organize the proof!
        \end{proof}
    \subsection{Pros and Cons}
        Pros:
        \begin{itemize}
            \item It learns safer paths by accounting for penalties from exploratory moves. This can be checked with cliff walking of \cite{Sutton2020}, with comparison to Q-Learning.
        \end{itemize}
        Cons:
        \begin{itemize}
            \item It learns suboptimal policies first, due to the behavior of searching for safe paths in the initial training process.
            \item The variance can be high, since in the initial training process the q values are very noisy and inaccurate, and the agent use the q value of it directly in the update.
                This can be dealt with the Expected Sarsa \ref{sec:ExpectedSarsa}.
        \end{itemize}